\subsection{不期而遇的邂逅}

也许是经过谈话后精神上感到疲惫，真昼比平时更早回了自己家，这个家里只剩下周和慧两人。\\

虽说如此，周和慧既没有共同点，又是今天刚认识，几乎算是陌生人，加上还有年龄差，所以并没有什么能分享的愉快话题，只是互相谦让着洗澡、刷牙等洗漱准备。

关于睡觉的地方，慧坚决拒绝睡床，而周则说睡沙发对小孩子的身体太残酷，最后双方妥协，在周的床边铺上客用被褥，弄得就像树来借宿时一样的状态。\\

两个并不亲近、说白了就是毫无关系、几乎算是陌生人的人睡在同一个房间里，有一种奇妙的紧张感，老实说周也不知道该怎么办。不过慧似乎也一样，他带着僵硬的笑容，郑重地道谢道「今天多谢你了。如果哥哥不嫌弃，请让我日后报答这份恩情」，然后背对这边钻进了被窝。\\

本以为是因为尴尬，但过了一会儿，听到了非常微弱的呜咽声，周这才恍然大悟，他原来是在心里整理关于自己家人的事情。\\

自己的父母对真昼做过的事，结合自己的常识去对照，他会有什么感受呢？自己今后该如何与家人相处呢？要思考的事情有很多，对自己的未来感到不安也是理所当然的。\\

他没有向周吐露这些，是因为周终究是个外人，而且他从本质上明白周是站在真昼那一边的。同时他也理解，这不是周能解决的问题。\\

这一切，都是必须由他自己去思考并解决的问题。\\

慧的父母对两个孩子都做出了残酷的事，周作为一个外人也不禁感到愤怒，但这并不是周能插嘴的问题，慧内心的整理只能由慧自己来完成。而且家庭之间的问题，只能由当事人自己商量解决。\\

现在周能做的，就只有静静地、不被他察觉地守望着他的内心纠葛。\\

听着那极其微弱的呜咽声，周轻轻咬紧嘴唇，闭上了眼睛。\\

\vspace{2\baselineskip}

第二天早上，看到起床后的慧，周觉得他的眼角似乎隐隐发红，但周对此什么也没提，只是把毛巾递给他，催促他去洗脸。

也许在周睡着之后他还在一直烦恼，但既然慧不开口，那就是慧自己决定的事，是他不想说出口的事。周打算尊重他的意愿，所以什么也没对他说。\\

「这次真的非常感谢。给你添了很多麻烦」\\

吃过早饭，到了人们开始活动的安全时间后才送他出门。在入口处，看着恭敬地低头致谢的慧，周和真昼都露出了微微的苦笑。\\

「住宿的事你不用在意。我还没冷酷到会在那种时候硬把你赶回去」

「这也算一方面，还有就是你应对了我这次突然的拜访。……给姐姐也添麻烦了」

「事到如今还说这些干什么。……我先声明，我比你想象的要冷静得多，所以请不要在意」

「……对不起」

「你向我道歉我也很为难。你对自己的行动没有后悔，对吧？」

「……是的」

「那我觉得这就够了。我的痛苦是我自己的问题。如果你也要面对你自己的问题，请优先处理那边吧」\\

这应该不是逞强，而是真昼经过一晚上的冷静整理后得出的真心话吧。仿佛在说「因为我有周在」，旁边的真昼用自己的手指轻轻碰了碰周的手指。周轻轻一笑，作为回应，也轻轻与她指尖相扣。\\

「慧君，回家之后才是真正的考验呢」

「是的。……在那之前，我得赶在被母亲发现之前赶紧回去。虽然我已经联络过了」\\

虽说姑且联络过说是要在朋友家过夜，但如果待得太久，也不能保证小夜那边不会联系对方的母亲。

现在是上午11点前。

因为给慧的衣服洗了又烘干，所以弄到了这个时间，但这难免会给人一种不像过夜，倒像赖着不走的印象。虽然还在勉强过关的范畴内，但如果再继续待下去，就有可能成为不在家被怀疑的因素。\\

「再次感谢你们昨天的照顾。答谢的事……」

「小孩子不用在意这些事。在把心思花在那上面之前，现在不是应该专注于你自己的事情上吗？」

「……你说得对。谢谢你」\\

还没解决问题的毕竟是慧这边。他带着有些苦涩的笑容鞠了一躬，说了句「那么我先告辞了」，然后转过身去，这时他却突然停下了脚步。\\

周正疑惑他是不是忘了什么东西，却注意到一辆车驶入了公寓的环岛，心想这应该是他想要避让车辆，想着想着，发现慧的样子有些不对劲。\\

慧回过头来，满脸写着惊愕。

旁边的真昼，触碰着的手指瞬间僵硬了。\\

那辆飒爽驶入的车在慧面前停下，车门滑开。

从后座上，走下了一位曾经见过的、留着柔顺茶色头发、表情充满生气的女性，她堂堂正正地走了下来。\\

「妈妈，为什么」

「我觉得情况不对劲，所以就多留心了一下。因为GPS显示你在奇怪的地方，我一瞬间还担心你是不是被绑架了呢」\\

「唉」小夜明显地叹了口气，她所说的话虽然是监视的那一套，但在义务教育期间，父母掌握孩子大致的去向也不能说是错的。只是这一次，这偏偏朝着对慧非常不利的方向发展了而已。\\

不用深思也能明白，慧应该是家境相当不错的少爷。为了防止被卷入奇怪的纠纷或事件，从预防的角度来看，事先设想要获取位置信息也是理所应当的。\\

本人之所以对此不知情，与其说是单纯地忘记了，倒不如说是极有可能根本不知道GPS的存在，加上父母在不需要的时候也不会插嘴，所以他才不知道。\\

「这地方还怪眼熟的。你就那么在意吗」

「……对不起，我瞒着你跑出来了」

「甚至不惜撒谎也要出来吗。真是的……要是这样的话，我是不是一开始就该告诉你才对，毕竟你也不小了」\\

「真是个让人头疼的孩子呢」小夜虽然苦笑着，但眼神中却没有任何厌恶的色彩，倒不如说流露出的是对自己的一种无奈。

那柔和的表情，和曾经展现给真昼的完全不同。无论从哪里看，都充满了无尽的慈爱，那是作为一位母亲的表情。

在旁边，传来微弱的握紧拳头的声音。\\

「不用听我也知道，慧想从那两个孩子那里问出些什么。也知道你想问我什么」\\

在掌握位置信息的那一刻，她大概就理解原因了吧。\\

「关于你的疑问，我不在现在回答」\\

并且她连慧得到了什么样的回答都理解了，抢先制止了慧想说的话。

慧原本打算回家后向小夜询问真相，同时也害怕彼此的关系会发生改变，但这全都是建立在小夜会回答的前提下的。如果小夜不回答，那么慧的决心和烦闷，就都变得毫无意义了。\\

「怎么会」

「我没说一直不回答。只是不该在现在这里说吧？」\\

直到这时，她的视线才第一次明确地投向了这边。准确地说，是投向了真昼。

看到那双与真昼不同、带着冷酷印象的眼睛径直望向真昼，周不由自主地挡在真昼面前，替她承受了那毫无温度的视线。

对方的眼睛眯得更细了，大概是因为周露出了严峻的表情吧。\\

「你的疑问，等回去冷静下来之后我再回答你。这样可以吗？」

「……嗯」\\

慧的目的是得知真相。既然得到了保证，他也没有理由反抗。

看到慧乖乖点头，小夜满意地眯起眼睛，轻轻叹了口气。\\

「很好。你先回去吧」

「哎，可是」

「我要向照顾了你一天的这两个人道谢之后再回去。好吗？」

「可是，那个」\\

慧担忧地抬头看向这边，特别是看向真昼。看到真昼那与昨天无法相比的、脸色苍白的模样，他也明白如果就这样让小夜和她说话会很不妙吧。\\

看到慧在担心真昼，小夜似乎觉得有些意外，守望着这边，但紧接着她又深深地吐了一口气，强烈得说是叹息也不为过。\\

「至少我不会对他们俩做什么。毕竟他们是慧的恩人。对吧？」

「……我可以相信你吧」

「我对你撒过谎吗？」

「从我懂事起，你不就在撒一个大谎吗」

「哎呀。……呵呵呵。是呢。你是对的，我是个大骗子。那就约定好吧。我不会加害这两个人。回去之后我会尽可能回答慧的问题。这样可以吗？……我没有打破过约定吧？」

「……嗯」\\

对于慧略带反抗的话语，小夜没有丝毫不悦，反而觉得很有趣、很高兴地笑着一带而过，这副模样任谁看了都觉得很有母亲的风范。

也许这类似于一种试探行为，但看到对方轻易地接受了，而且那大概也是他所熟知的平时母亲的模样，慧似乎安心了，他克制地点点头，小跑着来到小夜身边。\\

看到这一幕，真昼轻轻抓住了周衣服的袖子，显得很拘谨、无助地靠了过来。从上面俯视低着头的真昼，无法完全捕捉到她的表情，但可以想象，那至少是被归类在喜怒哀乐的哀之中的表情。\\

「那么，待会儿见。放心吧，不会变成你想象中的那种情况的」\\

小夜让慧坐进自己来时乘坐的那辆车，对司机简短地指示道「送他回家」。\\

透过后座的车窗，依然能看到慧担忧的脸，于是周为了让他安心而微笑着，表示真昼的事情就交给他。慧似乎理解到了周的意思，微微点了一下头。

车子很快从公寓的环岛消失了，在连车影都看不见的时候，小夜再次转向了这边。\\

这样面对面，再次感受到的是一种压迫感。\\

那种带刺的压迫感强烈到让真昼畏缩得发不出声，随着小夜深深地叹了口气，才多少有所缓和。\\

「真没想到，他竟然会跑到你这里来」\\

对真昼说出的话，比想象中还要平淡。\\

「那个孩子也真是让人头疼。从玲那里听说的时候我还以为怎么可能呢。这种行动力大概是随了玲吧」\\

从她的语气来看，玲应该就是慧的父亲。果然，那时玲可能已经隐约察觉到了。如果是通过GPS掌握了慧的位置，那他察觉到也不奇怪，在这个基础上，他先前可能是看穿了通话对象是周这边，所以才用了那种说法。\\

「算了。首先要道谢。感谢你们照顾那孩子。突然跑过来，给你们添麻烦了吧」\\

她带着似乎对这边既没有好感也没有恶意的态度，轻轻低下了头，但周丝毫没有放松警惕。只是，要是露出那种明显像是在炸毛一样的态度也很失礼，所以他始终在脸上贴着对外的笑容，缓缓摇了摇头。\\

「哪里谈得上麻烦。这终究是我们决定接受的事。所以，请不要在意」

「哎呀是吗。你还真是个认真的乖宝宝呢，虽然也想象过」

「……我就权当是称赞收下了」

「随你怎么想」\\

从想象过这句话就能听出，她确实仔细看过了报告书，大概连周的个人情况也都做过大致的调查吧。

对此周心里隐隐有些不快，但他没有表露出来。\\

「……那么，你呢？」\\

迟了一拍才被问到的真昼，明显地晃动了一下身体。

真昼说过已经不再把她当母亲看待了，但也许是唤醒了曾经的记忆，面对与小夜的直接对峙，她明显感到了强烈的抗拒感。只是，她并没有逃跑或者哭出来，更像是感到一种近乎恐惧的情绪。\\

被把话题抛过来，真昼明白不能一直保持沉默，她从周的影子里走出来，动作僵硬地低下了头。\\

「……好久不见，母亲。您看起来精神很好，那就再好不过」

「嗯，我很好。你看起来倒是一点都不好。嘛，这也是我和慧造成的，责怪你也不太好。不过那边的男朋友小哥，最好再把感情隐藏一下哦」

「……失礼了」\\

虽然没有说出口，但周的怪罪之意都写在了脸上，面对小夜那毫无掩饰的无奈的忠告，周再次紧紧闭上了嘴唇。\\

或许该说是幸好吧，小夜并没有特别生气的样子。或者说，她并没有把周他们当成对等的存在看待，这样说才更准确。她的眼神中带着一种「小孩子在搞什么名堂」的温度感。\\

从真昼的话中，周对她抱有相当难相处的印象，但实际交谈下来，虽然她的表情本身和遣词造句给人一种严厉的感觉，但态度本身并不是那么尖锐。要说这几分钟的感想，甚至可以说她是宽容的。\\

「一副好像有什么话想说的表情呢。有想说的话就说说看如何？」\\

虽然不想对慧那样承诺一定会回答，而且能感觉到根据内容和心情，周在意的事情可能一件都不会被揭露，但即便如此，她依然表现出了可以倾听的态度，这让周感到惊讶。

看起来并不像是在敌对，但也绝非抱有好意。面对态度平淡地给予许可的小夜，周虽然有些犹豫，但想到慧所追求的事，他还是开口了。\\

「我想向您请教一些事情。慧跟我说了许多，关于那些事」

「哎呀，真是贪心呢。你这个外人竟然要求我把连慧都还没告诉过的事情告诉你？要提出要求的话，也应该是那个孩子来提吧？那样我还更能接受一点」

「……我——」

「你想听吗？如果你想听，我可以告诉你，但对你来说，那绝对不是什么好事哦？」\\

小夜的说辞非常有道理，作为外人的周单方面抛出问题，她也没有义务回答。但是，如果是由作为亲生女儿、也是放弃抚养的受害者的真昼来问，那情况就不一样了。\\

她有知情的权利。而且，小夜也有诉说的意愿。\\

问题在于，身为当事人的真昼。\\

虽然想知道，但不想从小夜那里听到，不想待在她身边——这样的感情传递了过来，那份感情也表现在了行动上，她向后退了一步。

这大概是无意识的举动吧。但是，看着真昼那眼眸动摇、脸上失去血色的模样，如果继续让她面对小夜会很糟糕。更别提听小夜说话了。\\

「真昼。你想听吗？」

「我……」

「抱歉，我换个说法。想知道，但不用亲自听，这样可以吗？」\\

如果她希望，如果被允许的话，周打算成为真昼的耳目，代替真昼带回必要的信息。

似乎理解了周的意图，真昼带着毫无血色的脸，轻轻点了点头。这表示她同意了。

挺直腰背，再次将视线转向小夜，只见小夜正仔细地打量着这边，似乎觉得有些有趣。\\

「如果允许的话，我来代替真昼听。如果不行，那我就选择不听。您意下如何？」\\

如果小夜不说，那也就到此为止了。

无论真昼是在怎样的生长环境下活过来的，周都完整地爱着她，这并不会成为什么大问题。他们一直是在了解真昼遭受过放弃抚养这个事实的情况下共度的，也没有发生什么特别的问题。\\

也就是说，除了心理上的原因，并没有必要知道真昼被放弃抚养的理由，

接受了事实，周他们就会继续生活下去。至少，无论知不知道，没有被父母爱过的事实是无法改变的，而作为补偿，周会去爱真昼、珍惜真昼。真昼也接受了这样就好。\\

小夜会怎么接受这个提议呢——周毫不懈怠地保持着警戒，注视着她点燃了强烈意志的表情，结果小夜却很干脆地，显得轻松而大度地点了点头「也行吧」。\\

周没想到她会这么轻易地接受，不禁有些泄气，但小夜并不在意，而是将无比锋利的视线刺向真昼。\\

「你觉得，这样就好了，对吧？」

「……！」

「你真的，要全权委托给他吗？」\\

明明只是确认，但那语气中却包含了让真昼呼吸变得急促的压迫感，周不由自主地再次挡在真昼面前，遮挡住那视线。\\

「请别这样了。您看不见她的状况吗？」

「我看见了哦。看得很清楚呢。一软弱就闭门不出，简直和她父亲一模一样」\\

小夜向真昼抛出听起来只能算是讥讽的话语，但那声音和表情中却并没有恶意。对于这种矛盾的言行，周皱起了眉头，他把脸色苍白的真昼护在身后，保持着毅然的态度，将整个身体和脸笔直地转向小夜。\\

「到底是哪位把真昼变成这样的呢。能请您别施加不必要的压力了吗？」

「不要那么呲牙咧嘴的。如你所愿，我会陪你喝个茶的。这样你没意见了吧？」\\

意见当然是有的，但如果周和小夜一起离开这里，真昼就能从重压中解放出来。如果继续争论下去，交谈本身可能就会泡汤，与其错失好不容易抓到的良机，还是优先让真昼去避难比较好。\\

「之前听说那个笨蛋也和这个孩子一起喝过茶，还真是对自己宽容呢。一直移开视线又有什么用呢」

「您不也一样在移开视线吗？」

「哎呀。……呵呵，也许是这样呢？」\\

明明对她说了明显会让她感到不悦的话，小夜却愉快地笑着，说出模棱两可的话来观察这边的反应。

周评价她的视线像是在估价，但要说的话，这看起来似乎带有一种更加积极的意味。当然，这只是从周的视角来看。\\

「好吧，看在他的份上，我不会对你多说什么。你自己一个人好好想想吧」

「您……」\\

她退让得比想象中要快，从她的态度来看，看不出对真昼有什么执着。她似乎真的没有要攻击真昼的意思，不知是好是坏，她的兴趣反而转移到了周的身上。\\

「太过溺爱也是个问题哦。我懂你喜欢这孩子，但溺爱还是要在你能照顾得过来的范围内才行」

「我会带着觉悟和责任陪在她身边的」

「哦」\\

这简短冷淡的一句话，清晰地反映出她对周怎么对待真昼毫不在意的情感。\\

「真昼，你先去我家等着。没事的，好吗？」\\

明确表示了自己不在小夜的兴趣范围内，真昼虽然脸色依旧很差，但总比被迫面对要好得多。

希望能让她先在家里恢复平静，周为了让她安下心来，尽可能用温柔的微笑握住她的手，真昼轻微的颤抖慢慢地平息了下来。\\

面对等同于心理阴影般的存在并一直忍耐着，这说明真昼的精神已经能够保持相当的平衡了。只是，还没痊愈到可以面对面表达意见的程度而已。\\

「那我就把他借走咯。可以吧？」

「……好」\\

看到真昼轻轻点头，小夜露出了一丝类似淡淡的失落的色彩，这究竟是一种怎样的感情呢。\\

周知道就算问了也不会有回答，所以只用微弱的声音说了句「失陪了」，然后目送着转过身去的真昼消失在公寓内。\\

在真昼转过身的时候，小夜就立刻打电话安排了车，不到五分钟，另一辆车就停在了环岛。\\

「上来吧。我不会带你去什么奇怪的地方的」　

「这点我倒是不担心。……您打算去哪里？」

「我们集团旗下的一家料亭。在附近，能自由使用还能清场的地方，我只能想到这里了」\\

虽然早有预料，但小夜似乎拥有相当高的地位，周平时根本不会去的地方，她却能毫不费力地安排好，并且若无其事地对周说出来。

真昼曾说过唯独金钱方面不用发愁，所以周预想到她有着能够随便打发女儿也绰绰有余的收入，但她似乎是一位比想象中更有地位的女性。\\

「如果不愿意，在附近随便找个地方谈也可以，但你要是大声嚷嚷会引人注目的吧？」

「您是打算说一些会让我大声嚷嚷的事情吗」

「一扯到那孩子的事，你似乎就无法保持冷静呢。至少，你刚才像狗一样咬过来是事实吧？」\\

被这么一说，周无法反驳，但这一切的根源都是因为小夜和朝阳的所作所为，要求他消除愤怒也是强人所难。

似乎看穿了周的这种情绪，小夜丰润的嘴唇勾起了一抹妖异的弧线。\\

「毫不掩饰敌意可不是聪明的做法哦。就算是咬人，也最好选对时间和场合」

「……失礼了」

「算了，我也没心胸狭窄到要去计较一个小孩子的无礼，就当没看见吧」\\

说完，小夜毫不在意周的反应便坐进了车里，周虽然瞬间皱起了眉头，但也咽下了这口气，以免把感情表露得太多让对方如愿以偿，随后跟着她上了车。
